\subsection{Partial fixed-set Gamma results at weighted high temperature}
\label{ian:gamma-set-high-temperature}

Status: partial. Throughout the full stated weighted high-temperature
region, the Local alternative holds uniformly for degree-one
polynomials. It also holds quantitatively for polynomials with bounded
Fourier coefficient mass. Assuming (UAA), all degrees are covered
under an additional dimension-independent total-interaction budget.
The unrestricted row-sum-only Gamma assertion remains unresolved here.

Here (UAA) is the uniform theorem imported in \path{influence/high_temperature/01_conventions_imported_uniform_theorem.tex}, Section~\ref{ht:conventions-and-imported-uniform-theorem}.

Write
\[
 \mu(x)\propto
 \exp\left(\beta\sum_{|A|\ge2}J_Ax^A+\beta\sum_i h_ix_i\right),
 \qquad
 \omega=\beta\max_i\sum_{\substack{A\ni i\\|A|\ge2}}
                  (|A|-1)|J_A|\le\frac18.
\]
Parameters are finite, \(\beta\ge0\), and fields are arbitrary.

\begin{lemma}[A hereditary variance inequality]
\label{ian:weighted-dobrushin-poincare}
For every real function \(g\),
\[
 \Var_\mu g
 \le\frac1{1-\omega}
       \sum_i\E_\mu\Var(g\mid X_{-i})
 \le\frac1{1-\omega}\sum_iI_i^\mu(g).
 \tag{HTP}
\]
The same bound holds after every positive-probability pinning.
\end{lemma}

\begin{proof}
For \(i\ne j\), set
\[
 C_{ij}=\beta\sum_{A\ni i,j}|J_A|,\qquad C_{ii}=0.
\]
Changing boundary spin \(j\) changes the conditional law at \(i\)
in total variation by at most \(C_{ij}\): the conditional
plus-probability is \((1+\tanh u)/2\), whose derivative has absolute
value at most \(1/2\). The matrix \(C\) is symmetric and has row and
column sums at most \(\omega\).

Let \(P\) be random-scan heat bath. Couple its transitions from two
states differing only at \(j\), using the same update site and a
maximal coupling at that site. Updating \(j\) removes the disagreement;
updating \(i\ne j\) creates an additional disagreement with probability
at most \(C_{ij}\). Hence the expected Hamming distance is at most
\[
 q:=1-\frac{1-\omega}{n}.
\]
Joining adjacent-state couplings along Hamming paths gives contraction
by \(q\) in Hamming Wasserstein distance. Consequently
\[
 \operatorname{Lip}(Pg)\le q\,\operatorname{Lip}(g).
\]
For every nonconstant eigenfunction of the finite reversible kernel
\(P\), this implies that its eigenvalue has absolute value at most
\(q\). Thus the spectral gap is at least \((1-\omega)/n\).
The finite-state spectral variance inequality and
\[
 \langle g,(I-P)g\rangle_\mu
 =\frac1n\sum_i\E_\mu\Var(g\mid X_{-i})
\]
prove the first inequality in (HTP). The second follows from
\[
 \Var(g\mid X_{-i})
 =\Var(X_i\mid X_{-i})(D_ig)^2\le(D_ig)^2.
\]

Pinning reduces interaction supports. After merging equal residual
interactions, the triangle inequality shows that the weighted row
bound cannot increase. Newly created singleton terms are absorbed
into the fields. The same proof therefore applies after pinning.
\end{proof}

\begin{proposition}[Linear polynomials and bounded coefficient mass]
\label{ian:weighted-linear-aa}
Every globally bounded linear polynomial with
\(\Var_\mu f\ge v\) satisfies
\[
 \max_iI_i^\mu(f)\ge(1-\omega)^2v^2
 \ge\frac{49}{64}v^2.
\]
This assertion is uniform after pinning.

More generally, write the ordinary cube expansion as
\[
 f(x)=\sum_A\widehat f(A)x^A,\qquad
 M(f)=\sum_{A\ne\varnothing}|A|\,|\widehat f(A)|.
\]
If \(\Var_\mu f\ge v>0\), then
\[
 \max_iI_i^\mu(f)
 \ge\frac{(1-\omega)^2v^2}{M(f)^2}.
 \tag{HTC}
\]
In particular, if
\(\sum_A|\widehat f(A)|\le A_0d^q\), with \(q\ge0\), then
\[
 \max_iI_i^\mu(f)
 \ge\frac{49}{64A_0^2}(v/d)^{2q+2}.
\]
These assertions do not require (UAA).
\end{proposition}

\begin{proof}
Put \(I_*=\max_iI_i^\mu(f)\). Expanding one derivative and applying
Cauchy--Schwarz gives
\[
 \begin{split}
 I_i^\mu(f)
 &=\sum_{A\ni i}\widehat f(A)
       \E_\mu[(D_if)X^{A\setminus\{i\}}]\\
 &\le\sqrt{I_i^\mu(f)}
       \sum_{A\ni i}|\widehat f(A)|.
 \end{split}
\]
Therefore
\[
 \sum_iI_i^\mu(f)\le\sqrt{I_*}\,M(f).
\]
Combining this with (HTP) proves (HTC).

For \(f=a_0+\sum_i a_ix_i\), global boundedness implies
\[
 |a_0|+\sum_i|a_i|\le1,
\]
so \(M(f)\le1\). This proves the linear assertion. Restrictions
remain globally bounded and linear, so it is hereditary.

For the final assertion,
\(M(f)\le d\sum_A|\widehat f(A)|\le A_0d^{q+1}\);
use \(v\le1\) in (HTC).
\end{proof}

\begin{proposition}[All degrees with a total-interaction budget]
\label{ian:interaction-budget-aa}
Assume (UAA). Suppose, in addition, that
\[
 \beta\sum_{|A|\ge2}|J_A|\le B
\]
for a fixed dimension-independent \(B\).
Then every admissible polynomial satisfies
\[
 \max_iI_i^\mu(f)
 \ge\kappa_1e^{-26B}(v/d)^{12},
\]
where \(\kappa_1\) is defined in
Lemma~\ref{ian:block-resampling-aa}.
Consequently this subclass satisfies
\(\Gamma^{\mathrm{set}}\), and indeed
\(\Gamma^{\mathrm{rand}}\), through the Local alternative alone,
uniformly after pinning. No field bound or bias floor is needed.
\end{proposition}

\begin{proof}
Let
\[
 \nu(x)\propto e^{\beta\sum_i h_ix_i}
\]
be the corresponding product law, and put
\(W(x)=\beta\sum_{|A|\ge2}J_Ax^A\).
Since \(|W|\le B\),
\[
 e^{-2B}\le\frac{d\mu}{d\nu}\le e^{2B}.
\]
Thus
\[
 v\le\Var_\mu f
 \le\E_\mu(f-\E_\nu f)^2
 \le e^{2B}\Var_\nu f,
 \qquad
 I_i^\mu(f)\ge e^{-2B}I_i^\nu(f).
\]
Corollary~\ref{ian:two-spin-blocks}, applied to the arbitrary-bias
product \(\nu\), yields
\[
 \max_iI_i^\mu(f)
 \ge e^{-2B}\kappa_1
          \left(\frac{e^{-2B}v}{d}\right)^{12}
 =\kappa_1e^{-26B}(v/d)^{12}.
\]

Pinning cannot increase the total absolute mass of residual
interactions of size at least two: merge equal residual supports
using the triangle inequality and absorb singleton terms into the
fields. The same constant therefore works after pinning.
\end{proof}

\paragraph{Remaining scope.}
The weighted row bound does not supply a dimension-independent total
budget: arbitrarily many disjoint edges can satisfy the row bound
while their total interaction mass diverges. Nor does global
boundedness supply a dimension-independent polynomial bound on
\(M(f)\). For example, with \(H\) the \(n\times n\) Sylvester Hadamard matrix, \(f(x)=x^\top Hx/(2n^{3/2})\) satisfies \(|f|\le1/2\) on the whole cube because \(\|H\|_{\mathrm{op}}=\sqrt n\), while its off-diagonal coefficients have total absolute value \(n(n-1)/(2n^{3/2})\), of order \(\sqrt n\).

The star counterexample in Proposition~\ref{ian:forest-addressing}
has weighted row sum \(r\theta\), so it leaves the present region
as \(r\) grows. It therefore does not refute Gamma here.

The acceptance-simulation theorem in
\path{high_temperature/classical/05_arbitrary_degree.tex},
Theorem~\ref{thm:resolution-dynamics-arbitrary-degree}, remains valid
throughout the stated region. Proposition~\ref{ian:gamma-set-not-necessary}
shows why that simulation conclusion alone cannot establish fixed-set
Gamma. A proof or counterexample for arbitrary bounded low-degree
polynomials under the sole hypothesis \(\omega\le1/8\) remains needed.
